\chapter{Anexos}

\section{Código en Fortran con el método EMD}\label{EMD_code}

\begin{spacing}{1.}
\begin{lstlisting}
!============================================================================
!Module containing subroutines to perform the emd algorithm
!============================================================================
module emd_mod
  implicit none
  integer, parameter :: dp = selected_real_kind(15, 307)
  contains
!----------------------------------------------------------------------------
!Adds an element to an array
!----------------------------------------------------------------------------
  subroutine append(array, element)
    integer :: i, insize
    real(dp), intent(in) :: element
    real(dp), dimension(:), allocatable, intent(inout) :: array
    real(dp), dimension(:), allocatable :: outlist

    if(allocated(array)) then
      insize = size(array)
      allocate(outlist(insize+1))
      do i=1,insize          
        outlist(i) = array(i)
      end do
      outlist(insize+1) = element
      deallocate(array)
      call move_alloc(outlist, array) 
    else
      allocate(array(1))
      array(1) = element
    end if
  end subroutine append
!----------------------------------------------------------------------------
!Finds the local maxima and minima of an array and its positions in x-axis 
!----------------------------------------------------------------------------
  subroutine findLocalMaximaMinima(x, array, x_minima, minima, x_maxima, maxima)
    integer :: N, i, m, num_minima, num_maxima
    real(dp), intent(in) :: x(:), array(:)
    real(dp), dimension(:), intent(inout), allocatable :: x_minima, minima, &
       x_maxima, maxima
    real(dp), dimension(:), allocatable :: x_aux, y_aux 

    allocate(x_minima(0), minima(0), x_maxima(0), maxima(0))
    N = size(array)
    num_minima = 0
    num_maxima = 0
    do i=2, N-1
      !condition for local minima
      if((array(i-1)>array(i)).and.(array(i)<array(i+1))) then 
        call append(x_minima, x(i))
        call append(minima, array(i))
        num_minima = num_minima + 1
      !condition for local maxima
      else if((array(i-1)<array(i)).and.(array(i)>array(i+1))) then 
        call append(x_maxima, x(i))
        call append(maxima, array(i))
        num_maxima = num_maxima + 1
      end if
    end do
    
    !Boundary conditions to perform Cubic Splines
    !This adds 2 values to local maxima and minima: 4 in total
    !To get the exact number, the size of each must be reduced by  2
    if((num_minima .ne. 0).and.(num_maxima .ne. 0)) then
    !boundary conditions for minima
      m = size(x_minima)
      allocate(x_aux(m+2), y_aux(m+2))
      x_aux(1) = x(1)
      x_aux(2:m+1) = x_minima(1:m)
      x_aux(m+2) = x(N)
      y_aux(1) = minima(1)
      y_aux(2:m+1) = minima(1:m)
      y_aux(m+2) = minima(m)
      deallocate(x_minima, minima)
      allocate(x_minima(m+2), minima(m+2))
      x_minima = x_aux
      minima = y_aux 
      deallocate(x_aux, y_aux)

    !boundary conditions for maxima
      m = size(x_maxima)
      allocate(x_aux(m+2), y_aux(m+2))
      x_aux(1) = x(1)
      x_aux(2:m+1) = x_maxima(1:m)
      x_aux(m+2) = x(N)
      y_aux(1) = maxima(1)
      y_aux(2:m+1) = maxima(1:m)
      y_aux(m+2) = maxima(m)
      deallocate(x_maxima, maxima)
      allocate(x_maxima(m+2), maxima(m+2))
      x_maxima = x_aux
      maxima = y_aux 
      deallocate(x_aux, y_aux)
    end if
  end subroutine
!----------------------------------------------------------------------------
!Performs Cubic Splines using 'x', 'y' and 'x_splines' vectors as input and 
!returns 'y_splines' vector as output, interpolating the data in 'y'  
!----------------------------------------------------------------------------
  subroutine cubic_splines(x, y, x_splines, y_splines)
    real(dp), dimension(:), intent(in) :: x, y, x_splines
    real(dp), dimension(:), intent(inout) :: y_splines
    real(dp), dimension(:), allocatable :: b, h, c, d, l, dd, u
    integer :: i, j, k, Nnodes, N_splines
    real(dp) :: xinter, yinter

    Nnodes = size(x)
    N_splines = size(x_splines)
    allocate(b(Nnodes),h(Nnodes-1),c(Nnodes-1),d(Nnodes-1),l(Nnodes), &
      dd(Nnodes),u(Nnodes)) 

    !Tridiagonal matrix and vector of coefficients for the method
    !A(l, dd, u)x = b   
    do i = 1, Nnodes-1
      h(i) = x(i+1)-x(i)
      l(i+1) = 1.0_dp/h(i)
      u(i) = 1.0_dp/h(i)
      if(i>1) then 
        dd(i) = 1./h(i-1) + 1.0_dp/h(i)
        b(i) = y(i+1)/(h(i)*h(i)) + y(i)*(1./(h(i-1)*h(i-1))-1./(h(i)*h(i))) - &
          y(i-1)/(h(i-1)*h(i-1))
      end if
    end do !i
    dd(1) = 1.0_dp/h(1)
    dd(Nnodes) = 1.0_dp/h(Nnodes-1)
    dd(1:Nnodes) = 2.0_dp*dd(1:Nnodes)
    b(1) = (y(2)-y(1))/(h(1)*h(1))
    b(Nnodes) = (y(Nnodes)-y(Nnodes-1))/(h(Nnodes-1)*h(Nnodes-1))
    b(1:Nnodes) = 3.0_dp*b(1:Nnodes)

    !Solution of the linear system
    do i=2,Nnodes
      l(i) = l(i)/dd(i-1)
      dd(i) = dd(i) - l(i)*u(i-1)
      b(i) = b(i)-l(i)*b(i-1)
    end do !i
    b(Nnodes) = b(Nnodes)/dd(Nnodes)
    do i = Nnodes-1, 1, -1
      b(i) = (b(i) - u(i)*b(i+1))/dd(i)
    end do !i

    !Cubic Splines Interpolation
    do i = 1, Nnodes-1
      c(i) = 3.0_dp*(y(i+1)-y(i))/(h(i)*h(i))-2.0_dp*b(i)/h(i)-b(i+1)/h(i)
      d(i) = 2.0_dp*(y(i)-y(i+1))/(h(i)**3) + (b(i)+b(i+1))/(h(i)*h(i))
    end do 
    j = 1
    do i = 1, N_splines
      xinter = x_splines(i)
      do k = 1, Nnodes-1
        if((x(k)<xinter).and.(xinter < x(k+1))) j = k
      end do !k
      yinter = y(j) + b(j)*(xinter-x(j)) + c(j)*(xinter-x(j))*(xinter-x(j)) + &
        d(j)*(xinter-x(j))**3 
      y_splines(i) = yinter
    end do !i
    deallocate(h,b,c,d,l,dd,u)  
  end subroutine
!----------------------------------------------------------------------------
!Performs the main steps in the sifting process.
!It takes 'x' and 'y' as input vectors and returns a proto-IMF 'hk' and also
!'mean' and 'diff_extremas_zeros' in order to check if 'hk' is an IMF.
!----------------------------------------------------------------------------
  subroutine sifting_steps(x, y, hk, mean, diff_extremas_zeros)
    real(dp), dimension(:), intent(in) :: x, y
    real(dp), dimension(:), intent(inout) :: hk
    real(dp), dimension(:), allocatable :: x_minima, minima, x_maxima, maxima, &
      ys_low, ys_up, ys_mean
    integer, intent(inout) :: diff_extremas_zeros
    real(dp), intent(inout) :: mean
    integer :: N_file, i, number_zero_crossings

    N_file = size(x) 
    call findLocalMaximaMinima(x, y, x_minima, minima, x_maxima, maxima) 
    allocate(ys_low(N_file))
    call cubic_splines(x_minima, minima, x, ys_low)
    allocate(ys_up(N_file))
    call cubic_splines(x_maxima, maxima, x, ys_up)
    allocate(ys_mean(N_file))
    ys_mean = 0.5*(ys_low + ys_up)

    hk = y - ys_mean
    mean = 0.0_dp
    do i = 1, N_file
      mean = mean + hk(i)
    end do
    mean = mean/N_file

    number_zero_crossings = 0
    do i = 1, N_file - 1
      if(hk(i)*hk(i+1) < 0) number_zero_crossings = number_zero_crossings + 1
    end do

    diff_extremas_zeros = abs(size(maxima)+size(minima)-number_zero_crossings - 4)
    !We subtract 4 to correct the size of maxima and minima
    deallocate(x_minima, minima, x_maxima, maxima, ys_low, ys_up, ys_mean)
  end subroutine
!----------------------------------------------------------------------------
!Standard Deviation criterion to stop the sifting process
!----------------------------------------------------------------------------
  subroutine stopping_criterion1(hkm1, hk, SD)
    real(dp), dimension(:), intent(in) :: hkm1, hk
    real(dp), intent(inout) :: SD
    integer :: i, total_time
    real(dp), parameter :: eps=1e-8

    total_time = size(hk)
    SD = 0.0_dp
    do  i = 1, total_time
        SD = SD + abs(hk(i)-hkm1(i))*abs(hk(i)-hkm1(i))/(hkm1(i)*hkm1(i)+eps)
    end do
  end subroutine
!----------------------------------------------------------------------------
!Evaluates if the input parameters guarantee an IMF
!----------------------------------------------------------------------------
  subroutine is_imf(param1, param2, success_imf)
    !Difference between the number of extremas and zero crossings
    integer, intent(in) :: param1 
    real(dp), intent(in) :: param2 !mean of Proto-IMFs
    logical, intent(inout) :: success_imf
    real(dp), parameter :: eps_mean = 4e-2 !chosen mean tolerance

    if((param1 <= 1) .and. (abs(param2) <= eps_mean)) then 
      success_imf = .true.
    else
      success_imf = .false.
    end if
  end subroutine
!----------------------------------------------------------------------------
!Evaluates the number of extremas of the residue to stop the decomposition
!----------------------------------------------------------------------------
  subroutine residual_condition(x, r, num_residual_extremas)
    real(dp), dimension(:), intent(in) :: x, r
    integer, intent(inout) :: num_residual_extremas
    real(dp), dimension(:), allocatable :: rx_minima, r_minima, rx_maxima, r_maxima
    integer :: size_minima, size_maxima

    !Precision correction in the residual values. 
    !This avoids getting wrong IMFs.
    call findLocalMaximaMinima(x, anint(r*1e8)/1e8, rx_minima, r_minima, & 
      rx_maxima, r_maxima)
    size_minima = size(r_minima)
    size_maxima = size(r_maxima)
    !Size correction for minima and maxima
    if((size_minima.ne.0).and.(size_maxima.ne.0)) then 
      num_residual_extremas = size_minima + size_maxima - 4
    else
      num_residual_extremas = 0
    end if
    deallocate(rx_minima, r_minima, rx_maxima, r_maxima)
  end subroutine 
end module
!============================================================================
!Main program to perform EMD method on magnetic time series.
!============================================================================
program emd
  use emd_mod
  implicit none

  !MAX_IMFs~log2(N): in this case MAX_IMFs~ 16, 17: expected number of IMFs
  !MAX_PMFs: maximum number of iterations to perform the sifting process
  integer, parameter :: MAX_IMFs=16 + 2, MAX_PMFs=20
  real, parameter :: SD_threshold = 0.3_dp
  integer :: io1=1, io2=2
  integer :: N_file
  real(dp), dimension(:), allocatable :: x, y, hkm1, hk, rj, & 
    time, Br, Bn, Bt, B
  integer :: i, diff_extremas_zeros, num_residual_extremas, num_IMFs, it_sift
  character (len=80) :: nfout
  real(dp) :: mean, SD
  logical :: success_imf

  !size of time series data, name of output file with IMFs
  N_file = 62400 !PSP
  nfout = 'imfs_PSP_Br.dat'  
  !N_file = 115201 !SolO
  !nfout = 'imfs_SolO_B.dat'    
  nfout = trim(nfout)

  allocate(x(N_file), y(N_file), hk(N_file), hkm1(N_file), rj(N_file), &
    time(N_file), Br(N_file), Bt(N_file), Bn(N_file), B(N_file))

  !data reading
  open(unit=io1, file='pspr_20200927_0305_20200927_0515.txt', action='read')
  !open(unit=io1, file='solo_20201001_2120_20201002_0120.txt', action='read')
  do i = 1, N_file
    read(io1, *) time(i), Br(i), Bt(i), Bn(i), B(i)
  end do 
  close(io1)

  !variables on which the EMD method will be applied
  x(1:N_file) = time(1:N_file)/8.0_dp !sample rate = 8Hz
  y(1:N_file) = Br(1:N_file)

  open(unit=io2, file=nfout, status='unknown', action='write')
  
  !Residual condition: monotonous or constant behavior: at most 1 extrema
  num_residual_extremas = 2
  num_IMFs = 0
  rj = y !initial condition for residue
!---------------------------------EMD---------------------------------------
  do while((num_residual_extremas>1).and.(num_IMFs<MAX_IMFs))
    num_IMFs = num_IMFs + 1
    print*, 'IMF: ', num_IMFs
    hkm1 = rj

    !------------------------SIFTING----------------------------------------
    SD = 1.0_dp
    it_sift = 0
    do while((SD>SD_threshold).and.(it_sift<MAX_PMFs))
      it_sift = it_sift + 1
      hk = 0.0_dp
      call sifting_steps(x, hkm1, hk, mean, diff_extremas_zeros)
      call is_imf(diff_extremas_zeros, mean, success_imf)
      call stopping_criterion1(hkm1, hk, SD)
      !condition to stop sifting
      if(success_imf.eqv.(.true.)) exit
      hkm1 = hk
    end do !end of sifting process
    
    !report of sifting results 
    write(*, "(A5, I2, A10, I5, A8, E20.8, A7, L1, A12, F20.8)") &
      '#it: ', it_sift, '| param1: ', diff_extremas_zeros, &
        '| mean: ', mean, '| imf: ', success_imf, '| criterion: ', SD

    if(it_sift==MAX_PMFs) then 
      print*, 'Approximated sifting'
    else if (SD<0.3_dp) then 
      print*, 'SD condition achieved'
    end if

    !residue update
    rj = rj - hk
    write(io2, *) hk
    call residual_condition(x, rj, num_residual_extremas)
    print*, '|---------Extrema of residue---------|', num_residual_extremas
  end do !end of EMD

  write(io2, *) rj

  print*, 'Extrema of residue: ', num_residual_extremas
  print*, 'Number of IMFS: ', num_IMFs

  deallocate(x, y, hkm1, hk, rj, time, Br, Bn, Bt, B)
  close(io2)
end program
!============================================================================
\end{lstlisting}
\end{spacing}

\section{Código en Python para realizar el barrido sobre las ventanas de tiempo y el número de IMFs}\label{xcorr_scan_code}

\begin{spacing}{1.}
\begin{minted}{python}
#=============================================================================
#Program to perform the cross-correlation scan on the time windows and on 
#the number of IMFS
#=============================================================================
#libraries
import numpy as np
import emd as emd2
import numba
from numba import jit
import seaborn as sns; sns.set_theme()

#data
psp_data = np.loadtxt('pspr_20200927_0305_20200927_0515.txt')
psp_data = np.transpose(psp_data)
psp_data[0] = psp_data[0]/8.
solo_data = np.loadtxt('solo_20201001_2120_20201002_0120.txt')
solo_data = np.transpose(solo_data)
solo_data[0] = solo_data[0]/8.
#-----------------------------------------------------------------------------
#Functions
#-----------------------------------------------------------------------------
#Cross-Correlation function
@numba.jit
def cross_corr(x, y):
  n = len(x)
  mean_x = np.mean(x)
  mean_y = np.mean(y)
  std_x = np.std(x)
  std_y = np.std(y)

  xcorr = []
  for k in range(n):
    suma = 0.
    for i in range(n-k):
      suma += (x[i]-mean_x)*(y[i+k]-mean_y)
    xcorr.append(suma)
  xcorr = np.array(xcorr)/(n*std_x*std_y)
  return xcorr

#Function to get the partitions of PSP and SolO (portion of the signals) 
#of the correlation values from the maximum to the minimum
def get_jixcorrs(W):
  length = np.size(W)
  aux = np.reshape(W, (1, length))
  aux2 = np.sort(aux, kind='quicksort')
  aux3 = [np.where(W == aux2[0][i]) for i in range(length)]
  indexs = []
  for k in range(length):
    i, = aux3[k][1]
    j, = aux3[k][0]
    indexs.append([j, i, W[j][i]])
  indexs = indexs[::-1]
  return indexs
#Funtion to perform the cross-correlation scan on the number of IMFS
def params_xcorr(field_component, time_min, delay_min):
  #EMD parameters
  config = emd2.sift.get_config('sift')
  config['imf_opts/sd_thresh'] = 0.3

  #parameters
  T1 = 62400 #130 min PSP
  T2 = 115201 #240 min SolO
  time_width = time_min*60*8
  delay = delay_min*60*8 #time to change from one partition to another 
  n_inter1 = int(np.floor((T1-time_width)/delay)) + 1
  n_inter2 = int(np.floor((T2-time_width)/delay)) + 1
  Tp1 = time_width + delay*(n_inter1-1) #<=T1
  Tp2 = time_width + delay*(n_inter2-1) #<=T2

  #data
  x1 = psp_data[field_component][:Tp1]
  x2 = solo_data[field_component][:Tp2]
  
  #EMD
  imf1 = emd2.sift.sift(x1) #Residue included
  imf2 = emd2.sift.sift(x2)
  num_imfs1 = imf1.shape[1]-1
  num_imfs2 = imf2.shape[1]-1
  start_index_imfs1 = [i for i in range(num_imfs1)]
  start_index_imfs2 = [i for i in range(num_imfs2)]

  #SCAN
  params = []
  print('Time window: {} minutes'.format(time_min))
  sum_imfs = min(start_index_imfs1[-1], start_index_imfs2[-1])
  for k in range(sum_imfs):
    print('# imfs added : PSP({}), SolO({})'.format(num_imfs1-k, num_imfs2-k))

    #Elimination of high frequencies
    y1 = np.sum(imf1[:, start_index_imfs1[k] : ], axis=1)
    y2 = np.sum(imf2[:, start_index_imfs2[k] : ], axis=1)

    #cross-correlation between partitions over the time series
    #scan for a particular time window
    xcorrs = []
    max_lags = []
    for i in range(n_inter1):
      start1 = i*delay
      end1 = time_width + i*delay
      for j in range(n_inter2):
        start2 = j*delay
        end2 = time_width + j*delay
        aux = cross_corr(y1[start1:end1], y2[start2:end2])
        index, = np.where(aux == np.max(aux))
        xcorrs.append(np.max(aux))
        max_lags.append(index[0])

    xcorr_map = np.reshape(xcorrs, (n_inter1, n_inter2))
    xcorr_lags = np.reshape(max_lags, (n_inter1, n_inter2))
    max_vals = get_jixcorrs(xcorr_map)

    #map of maximum correlations
    ticks_step = 8
    xlabels0 = np.arange(time_min+1, time_min+delay_min*(n_inter2-1), delay_min, dtype=int)
    ylabels0 = np.arange(time_min+1, time_min+delay_min*(n_inter1-1), delay_min, dtype=int)
    xlabels1 = np.char.mod('%d', xlabels0)
    ylabels1 = np.char.mod('%d', ylabels0)
    xticks = np.arange(0, n_inter2, 1)
    yticks = np.arange(0, n_inter1, 1)

    fig, ax = plt.subplots(figsize=(width, heigth))
    ax = sns.heatmap(xcorr_map, vmin=0, vmax=1, cmap='inferno', 
      cbar_kws={'label': "correlation", 'pad': 0.005}, annot_kws={"size": 17})
    plt.text(0.35, 0.9, "Maximum = {:.3f}".format(np.max(xcorr_map)), fontsize=17,
       transform=plt.gcf().transFigure)
    ax.invert_yaxis()
    ax.set_xlabel(r"$t(SolO) \; [min]$", fontsize=17)
    ax.set_ylabel(r'$t(PSP) \; [min]$', fontsize=17)  

    ax.set_yticks(yticks[1::ticks_step])
    ax.set_yticklabels(ylabels1[::ticks_step])
    ax.set_xticks(xticks[1::ticks_step])
    ax.set_xticklabels(xlabels1[::ticks_step])
    plt.yticks(rotation=0)
    #save map of maximum correlations in img folder
    plt.savefig('img/fig-'+ str(num_imfs1 - start_index_imfs1[k]) + '&' + 
      str(num_imfs2 - start_index_imfs2[k]) + '_' + str(time_min) + '&'+ 
      str(delay_min) + '.png')
    plt.close(fig)
    
    params.append([[time_min, delay_min, num_imfs1 - start_index_imfs1[k], 
      num_imfs2 - start_index_imfs2[k]], max_vals, xcorr_map, xcorr_lags])
  return params

#Numba function test: needed for it to perform better
x1 = psp_data[1]
x2 = solo_data[1]
test = cross_corr(x1, x2)
#-----------------------------------------------------------------------------
#Main part of the program: It performs the cross-correlation scan 
#on the time windows and on the number of IMFS for some time windows.
#-----------------------------------------------------------------------------
time_mins = np.arange(10, 90, 2)
delay_min = 2
params_time = []

for i in time_mins:
    params_time.append(params_xcorr(1, i, delay_min))
#=============================================================================
\end{minted}
\end{spacing}